% ADD THIS HEADER TO ALL NEW CHAPTER FILES FOR SUBFILES SUPPORT

% Allow independent compilation of this section for efficiency
\documentclass[../CLthesis.tex]{subfiles}

% Add the graphics path for subfiles support
\graphicspath{{\subfix{../images/}}}

% END OF SUBFILES HEADER

%%%%%%%%%%%%%%%%%%%%%%%%%%%%%%%%%%%%%%%%%%%%%%%%%%%%%%%%%%%%%%%%
% START OF DOCUMENT: Every chapter can be compiled separately
\begin{document}

\phantomsection% to get the hyperlinks (and bookmarks in PDF) right for index, list of files, bibliography, etc.
\addcontentsline{toc}{chapter}{Abstract}

\begin{abstract}
    \begin{otherlanguage}{english}% Switch to English, in case the main language is not English

        \section*{Abstract}

        Computational hate speech research has made substantial progress on detection, yet it remains concentrated on English-language data, explicit lexical abuse, and binary classification --- a combination that leaves implicit, cross-lingual religious hostility structurally underexplored. This thesis investigates anti-Christian hate speech across Chinese, Japanese, and English online discourse, asking not only whether hostility can be detected, but \emph{who} is targeted, \emph{how} attacks are rhetorically constructed, and \emph{why} different communities sustain them.

        To answer these questions, the thesis assembles a trilingual corpus of 63,265 texts and develops a modular \emph{who--how--why} analytical pipeline. Low-resource Chinese and Japanese subsets are expanded through LLM-assisted back-translation augmentation, raising hate speech detection F1 from 0.172 to 0.635 (Chinese) and from 0.294 to 0.579 (Japanese). Unsupervised BERTopic modelling with cross-lingual centroid alignment identifies 47 semantically coherent topics without imposing a predefined English-centric taxonomy. A three-layer extraction architecture --- combining spaCy NER, a domain gazetteer, and a large language model fallback --- addresses the zero-subject constructions that make standard dependency parsing unreliable in CJK text. Attack frames are classified using a ten-category taxonomy, and moral motivations are projected onto classical Moral Foundation axes extended with intergroup threat dimensions.

        The results reveal two structural regularities that prior binary detection cannot surface. First, attack frame selection is systematically governed by the social category of the target: clerical role-targets attract the Sexual Abuse Frame across all three languages, institutionally-embedded targets attract Political Interference and Economic Exploitation, and collective identity labels attract Rhetorical Attack and Moral Accusation --- a genre-like target-to-frame regularity that holds cross-linguistically. Second, English and East Asian communities inhabit fundamentally different moral architectures: English discourse aligns negatively across all five classical Moral Foundation axes simultaneously, constituting a broad ethical condemnation from within a shared normative framework, while Chinese and Japanese discourse aligns negatively on Symbolic and Realistic Threat axes rather than on classical foundations, framing Christianity as a danger to cultural identity and social cohesion rather than as a moral transgressor.

        These findings are supported by a trilingual annotated dataset and an open-source pipeline released at \url{https://github.com/notsofun/Master_Thesis}, providing reusable resources for cross-lingual hate speech research in low-resource settings.

    \end{otherlanguage}
    %  A German Abstract is optional if there is an English one.
    %  \newpage \thispagestyle{empty}%if both languages need more than one page

    % \begin{otherlanguage}{ngerman}% Switch to German for this part

    %    \section*{Zusammenfassung}
    %    Und hier sollte die Zusammenfassung auf Deutsch erscheinen. Kann auch auf separater Seite erscheinen.

    % \end{otherlanguage}

\end{abstract}

\newpage

\end{document}
